%% =====================================================================
%%  T-21-04  The Frozen Crowd
%%  -------------------------------------------------------------------
%%  One idea per file. Core is mandatory. Slots in canonical order.
%%  Max 700 words. No layout commands. No filler. New source material
%%  for this entry goes in sources/B4/T-21-04.tex -- never by
%%  rewriting this file (docs/RULES.md R0.1).
%%  =====================================================================
%% @kind: topic
%% @id: T-21-04
%% @title: The Frozen Crowd
%% @subtitle: Bystander diffusion --- everyone looks at everyone, and nobody moves
%% @chapter: c21
%% @order: 40
%% @status: stable
%% @sources: B4
%% @updated: 2026-09-19

\topic{T-21-04}{The Frozen Crowd}{Bystander diffusion --- everyone looks at everyone, and nobody moves}

\begin{core}
In a crowd, the check everyone runs is what everyone else is doing --- and when everyone is uncertain, everyone reads everyone else's uncertainty as evidence that nothing is wrong. The result is the frozen crowd: thirty-eight neighbours watched a killing through lit windows and not one called the police in time. The mechanism --- pluralistic ignorance and responsibility divided by headcount --- is a lever for whoever benefits from inaction, and a countermeasure for whoever refuses it.
\end{core}
\begin{mechanism}
B4's laboratory source: the Genovese case --- attacked three times over thirty-five minutes, thirty-eight witnesses, one call, after she was dead. The mechanism has two gears. Pluralistic ignorance: in ambiguous situations each person checks the others; calm-looking faces all around are misread as \emph{everything is fine}, when each face is privately alarmed and publicly mirroring the same misread. Diffusion of responsibility: once it is clear something is wrong, \emph{someone} should act --- and the more someones, the less any single one feels the weight; help-seeking addressed to a group is help-seeking addressed to nobody. The laboratory result is the dose-response curve: a person in distress gets help quickly alone with one witness, slowly or never with several. The operator's version needs no malice: the aggressor who strikes in a crowd benefits automatically; the meeting where everyone waits for someone to name the obvious problem runs the same code. The countermeasure is pointed, and B4's field tests confirm it: dissolve the diffusion by making one person individually responsible --- a name, a finger, a direct request with eyes on.
\end{mechanism}
\begin{conditions}
Strongest with ambiguity (no one is sure of the script) and numbers (three or more witnesses is enough). Urban and online environments maximise both. It weakens when the need is unambiguous and the required act is cheap and specific, and it is defeated outright by individuation: a crowd addressed by name is a queue of individuals.
\end{conditions}
\begin{application}
  \item Need help in a crowd? Pick one person, describe them to themselves --- \emph{you, in the grey coat: call an ambulance} --- and hold their eyes.
  \item In meetings, the obvious problem nobody names: name it, then assign it --- \emph{X, can you own checking this by Friday?} General appeals produce the frozen crowd's office variant.
  \item When you feel the freeze yourself --- the corridor incident, the suspicious email --- act as if you are the only witness, because behaviourally you are.
  \item As a bystander who is frozen, note it aloud: \emph{I'm assuming someone has already called} is the freeze talking. Say the uncertain thing first.
  \item Aggressors know this map. If you must confront or de-escalate in public, do it where diffusion works for you --- and refuse to be the crowd for anyone else's aggression.
\end{application}
\begin{feedback}
  \item Working: a named person acts within seconds where a room full of people sat still.
  \item Working: naming the ambiguity in a meeting breaks the silence --- two others were waiting for permission.
  \item Failing: your general appeal to the group produced three follow-up conversations and no action. Diffusion was not dissolved.
  \item Failing: you left a scene telling yourself someone else had it. That sentence is the mechanism, speaking in your voice.
\end{feedback}
\begin{failure}
The countermeasure has its own cost: individuated requests concentrate real risk on the person named, and asking a frozen crowd for help can recruit its hostility instead --- crowds defend their inaction. The entry also explains atrocities' audiences as well as emergencies', and the honest limit is that knowledge of the mechanism does not immunise the knower: the freeze is felt from inside as caution, not cowardice.
\end{failure}
\begin{countermeasures}
  \item Train the pointed request until it is automatic. In crisis you will not design the sentence; you will run whichever one is loaded.
  \item Refuse to be an extra witness: if you would want help, give it --- the first mover re-prices the crowd's normals instantly.
  \item In organisations, ban the orphan task: anything addressed to \emph{the team} needs a name attached before the meeting ends.
  \item Notice manufactured ambiguity: bad news framed so no one can tell whether acting is overreacting is the freeze being farmed (T-15-02).
\end{countermeasures}

\sources{B4}
\seesources
\seealso{T-15-01, T-15-02, T-13-01}
